\chapter{프라이머리 분해와 순환 분해}

\chapterintro{이 장에서는 최소다항식의 인수분해 정보를 이용해 선형연산자를 큰 블록으로 나누고, 다시 순환 블록으로 정밀하게 분해하겠습니다. 먼저 소멸 아이디얼(annihilator ideal)과 순환 부분공간(cyclic subspace)을 정리한 뒤, 프라이머리 분해(primary decomposition) 정리와 순환 분해(cyclic decomposition) 정리를 연결하겠습니다. 이 흐름을 이해하면 다음 장의 Jordan 표준형과 유리 표준형으로 자연스럽게 이어집니다.}
\chaptergoals{소멸 아이디얼과 순환 부분공간의 관계를 엄밀히 설명할 수 있다.}{최소다항식 인수분해로부터 프라이머리 분해를 구성하고 증명할 수 있다.}{순환 분해 정리를 통해 companion block 표현의 의미를 해석할 수 있다.}

\sectiontemplate{소멸 아이디얼과 순환 부분공간}{순환 분해의 출발점은 ``한 벡터가 만드는 궤도''를 다항식 언어로 읽는 것입니다. 소멸 아이디얼을 정확히 잡아 두면 차원 계산과 기저 구성이 동시에 정리됩니다.}{\item 소멸 아이디얼의 정의를 쓸 수 있습니다.\item 순환 부분공간의 기저를 구성할 수 있습니다.\item 서로소 다항식에 대한 kernel 분해 공식을 적용할 수 있습니다.}{\item 유클리드 알고리즘\item 최소다항식\item 불변부분공간}

\begin{definition}
$T\in\operatorname{End}(V)$, $S\subseteq V$에 대하여
\[
\operatorname{Ann}_T(S)=\{f(t)\in F[t]\mid f(T)s=0\ \text{for all }s\in S\}
\]
를 $S$의 $T$-소멸 아이디얼(annihilator ideal)이라 한다. 특히 $v\in V$에 대해 $\operatorname{Ann}_T(v)=\operatorname{Ann}_T(\{v\})$라 쓴다.
\end{definition}

\begin{definition}
$v\in V$에 대하여
\[
Z(v;T)=\operatorname{span}\{v,Tv,T^2v,\dots\}
\]
를 $v$가 생성하는 $T$-순환 부분공간(cyclic subspace)이라 한다. $Z(v;T)=V$이면 $v$를 순환 벡터(cyclic vector)라 한다.
\end{definition}

\begin{theorem}
$0\ne v\in V$라 하자. 그러면 다음이 성립한다.
\begin{enumerate}[label=\arabic*.]
\item $\operatorname{Ann}_T(v)$는 $F[t]$의 영이 아닌 아이디얼이며, 유일한 단위다항식 $m_{T,v}$가 존재하여
\[
\operatorname{Ann}_T(v)=(m_{T,v})
\]
이다.
\item $d=\deg m_{T,v}$라 두면
\[
\{v,Tv,\dots,T^{d-1}v\}
\]
는 $Z(v;T)$의 기저이다.
\end{enumerate}
\end{theorem}

\paragraph{증명 전략.}
유한차원성으로 먼저 비자명한 소멸다항식의 존재를 확보하고, $F[t]$가 PID라는 사실로 생성원을 단일 다항식으로 고정합니다. 그다음 다항식 나눗셈을 이용해 생성성과 독립성을 분리하여 순환 기저를 완성합니다.

\begin{proof}
$\dim V=n$이라 하자. 벡터열
\[
v,Tv,\dots,T^n v
\]
는 $n+1$개이므로 일차종속이다. 따라서 계수들이 모두 0이 아닌 것은 아닌 $c_0,\dots,c_n\in F$가 존재하여
\[
c_0v+c_1Tv+\cdots+c_nT^n v=0.
\]
즉 $f(t)=c_0+c_1t+\cdots+c_nt^n$에 대해 $f(T)v=0$이므로 $\operatorname{Ann}_T(v)\ne\{0\}$이다. 또한 $f(T)v=0$, $g(T)v=0$이면 $(f+g)(T)v=0$, 임의의 $h\in F[t]$에 대해 $(hf)(T)v=h(T)f(T)v=0$이므로 $\operatorname{Ann}_T(v)$는 아이디얼이다. $F[t]$는 PID이므로
\[
\operatorname{Ann}_T(v)=(m_{T,v})
\]
가 되는 단위다항식 $m_{T,v}$가 유일하게 존재한다.

이제 $d=\deg m_{T,v}$라 두자. 임의의 $k\ge d$에 대해 다항식 나눗셈으로
\[
t^k=q_k(t)m_{T,v}(t)+r_k(t),\qquad \deg r_k<d
\]
가 된다. $v$에 대입하면
\[
T^k v=q_k(T)m_{T,v}(T)v+r_k(T)v=r_k(T)v\in\operatorname{span}\{v,\dots,T^{d-1}v\}.
\]
따라서 $Z(v;T)$는 $\{v,\dots,T^{d-1}v\}$로 생성된다.

독립성을 보이자. 만약
\[
\alpha_0 v+\alpha_1 Tv+\cdots+\alpha_{d-1}T^{d-1}v=0
\]
가 비자명하다면 $r(t)=\alpha_0+\alpha_1 t+\cdots+\alpha_{d-1}t^{d-1}\ne0$이고 $r(T)v=0$이다. 즉 $r\in\operatorname{Ann}_T(v)=(m_{T,v})$이므로 $m_{T,v}\mid r$여야 한다. 그런데 $\deg r<\deg m_{T,v}$이므로 모순이다. 따라서 계수는 모두 0이고 집합은 일차독립이다. 결론적으로 $\{v,\dots,T^{d-1}v\}$는 $Z(v;T)$의 기저이다.
\end{proof}

\paragraph{의미 해석.}
한 벡터의 궤도는 ``몇 번까지 독립이 유지되는가''로 정확히 측정됩니다. 그 경계 차수가 바로 $m_{T,v}$의 차수이며, 이 수가 순환 블록의 크기를 결정합니다.

\begin{theorem}
$f,g\in F[t]$가 서로소라 하자. 그러면
\[
\ker\big((fg)(T)\big)=\ker f(T)\oplus\ker g(T)
\]
가 성립한다.
\end{theorem}

\paragraph{증명 전략.}
서로소 조건에서 나오는 Bézout 항등식을 $T$에 대입해 임의의 벡터를 두 kernel 성분으로 분해합니다. 그다음 교집합 벡터에 동일한 항등식을 적용해 자명교집합을 확인합니다.

\begin{proof}
$\gcd(f,g)=1$이므로 어떤 $a,b\in F[t]$가 존재하여
\[
a(t)f(t)+b(t)g(t)=1
\]
이다. 임의의 $x\in\ker((fg)(T))$를 잡으면
\[
x=(af+bg)(T)x=a(T)f(T)x+b(T)g(T)x.
\]
첫 항에 $g(T)$를 작용시키면
\[
g(T)a(T)f(T)x=a(T)(fg)(T)x=0
\]
이므로 $a(T)f(T)x\in\ker g(T)$이다. 같은 방식으로 $b(T)g(T)x\in\ker f(T)$이다. 따라서
\[
\ker((fg)(T))\subseteq \ker f(T)+\ker g(T).
\]
반대 포함은 자명하다. 실제로 $u\in\ker f(T)$, $w\in\ker g(T)$이면
\[
(fg)(T)(u+w)=f(T)g(T)u+f(T)g(T)w=0+0=0.
\]
이제 교집합을 보자. $y\in\ker f(T)\cap\ker g(T)$이면
\[
y=(af+bg)(T)y=a(T)f(T)y+b(T)g(T)y=0.
\]
따라서 교집합은 $\{0\}$이고 직합이 성립한다.
\end{proof}

\paragraph{의미 해석.}
서로소 다항식은 작용을 충돌 없이 분리합니다. 이 정리는 이후 프라이머리 성분을 ``겹침 없이'' 붙여 만드는 핵심 기계로 사용됩니다.

\begin{example}
$V=F^3$에서
\[
T=\begin{bmatrix}
0&1&0\\
0&0&1\\
0&0&0
\end{bmatrix}
\]
라 두고 $v=e_3$를 택하자. 그러면
\[
Tv=e_2,\qquad T^2v=e_1,\qquad T^3v=0
\]
이므로 $\operatorname{Ann}_T(v)=(t^3)$, $m_{T,v}=t^3$, $\{v,Tv,T^2v\}$는 $V$의 기저이다. 따라서 $v$는 순환 벡터이다.
\end{example}

\subsection*{주의}
\begin{warning}
이 절의 소멸 아이디얼은 쌍대공간에서 쓰는 annihilator와 대상이 다릅니다. 여기서는 ``다항식이 벡터를 소멸시키는 조건''을 다룹니다.
\end{warning}
\begin{warning}
$m_{T,v}$는 $\operatorname{Ann}_T(v)$의 임의 원소가 아니라 단위 생성원입니다. 최고차항 계수를 1로 고정하지 않으면 유일성이 깨집니다.
\end{warning}

\subsection*{자가진단퀴즈}
\begin{enumerate}[label=\arabic*.]
\item $\dim V=5$이고 $\deg m_{T,v}=3$이면 $\dim Z(v;T)$를 구하고 이유를 설명하십시오.
\item $f,g$가 서로소일 때 $\ker((fg)(T))=\ker f(T)+\ker g(T)$가 되는 핵심 항등식을 쓰십시오.
\item $v\ne0$에 대해 $\operatorname{Ann}_T(v)\ne\{0\}$가 되는 이유를 유한차원성만으로 설명하십시오.
\end{enumerate}

\sectiontemplate{프라이머리 분해 정리}{최소다항식을 기약다항식 거듭제곱으로 분해하면, 연산자의 작용은 서로 독립적인 프라이머리 성분(primary component)으로 갈라집니다. 이 절에서는 분해의 존재와 구성 연산자를 엄밀히 확인하겠습니다.}{\item 프라이머리 성분을 정의하고 불변성을 확인할 수 있습니다.\item Primary decomposition 정리를 완전하게 증명할 수 있습니다.\item 성분 제한 연산자의 최소다항식을 계산할 수 있습니다.}{\item 최소다항식 인수분해\item Bézout 항등식\item 직합 분해}

\begin{definition}
$T\in\operatorname{End}(V)$의 최소다항식이
\[
m_T(t)=p_1(t)^{e_1}\cdots p_r(t)^{e_r}
\]
로 분해된다고 하자. 여기서 $p_1,\dots,p_r$는 서로 다른 단위 기약다항식이다. 이때
\[
V_i=\ker p_i(T)^{e_i}\qquad (1\le i\le r)
\]
를 $T$의 $i$번째 프라이머리 성분(primary component)이라 한다.
\end{definition}

\begin{theorem}[Primary decomposition 정리]
위 표기에서
\[
V=V_1\oplus\cdots\oplus V_r
\]
가 성립한다.
\end{theorem}

\paragraph{증명 전략.}
각 인수 하나를 제외한 곱 $q_i=m_T/p_i^{e_i}$를 만들고, $\sum b_iq_i=1$ 형태의 다항식 항등식을 구성합니다. 이를 $T$에 대입해 성분 사영 연산자 $E_i$를 만든 뒤, 합이 전체공간이며 교집합이 자명함을 순서대로 확인합니다.

\begin{proof}
각 $i$에 대해
\[
q_i(t)=\frac{m_T(t)}{p_i(t)^{e_i}}
\]
라 두자. $q_1,\dots,q_r$가 생성하는 아이디얼은 $F[t]$ 전체이므로(즉 $\gcd(q_1,\dots,q_r)=1$) 어떤 $b_1,\dots,b_r\in F[t]$가 존재하여
\[
b_1q_1+\cdots+b_rq_r=1
\]
이다. 이제
\[
E_i=b_i(T)q_i(T)\in\operatorname{End}(V)
\]
로 둔다. 위 항등식에 $T$를 대입하면
\[
E_1+\cdots+E_r=I.
\]

먼저 $E_i(V)\subseteq V_i$를 보이자. 임의의 $x\in V$에 대해
\[
p_i(T)^{e_i}E_i x=b_i(T)m_T(T)x=0
\]
이므로 $E_i x\in\ker p_i(T)^{e_i}=V_i$이다.

또한 $j\ne i$이고 $x\in V_i$이면 $q_j$는 $p_i^{e_i}$를 인수로 가지므로
\[
q_j(T)x=0,\qquad E_jx=b_j(T)q_j(T)x=0.
\]
따라서 $x\in V_i$에 대해
\[
x=(E_1+\cdots+E_r)x=E_ix
\]
가 된다.

이제 합공간을 보자. 임의의 $x\in V$에 대해
\[
x=\sum_{i=1}^r E_i x,\qquad E_i x\in V_i
\]
이므로 $V=\sum_i V_i$이다.

직합성 확인: $\sum_{i=1}^r x_i=0$이고 $x_i\in V_i$라 하자. 고정된 $k$에 대해 $E_k$를 작용시키면
\[
0=E_k\Big(\sum_{i=1}^r x_i\Big)=\sum_{i=1}^r E_kx_i=E_kx_k=x_k.
\]
모든 $k$에 대해 $x_k=0$이므로 합은 직합이다. 따라서
\[
V=V_1\oplus\cdots\oplus V_r.
\]
\end{proof}

\paragraph{의미 해석.}
Primary decomposition은 ``최소다항식 인수 하나당 독립된 역학 블록 하나''라는 구조를 보장합니다. 이후 계산은 각 블록에서 따로 수행한 뒤 직합으로 합치면 됩니다.

\begin{theorem}
$T_i=T|_{V_i}$라 하자. 그러면
\[
m_{T_i}(t)=p_i(t)^{e_i}\qquad (1\le i\le r)
\]
이며, 동시에
\[
m_T=\operatorname{lcm}\big(m_{T_1},\dots,m_{T_r}\big)
\]
가 성립한다.
\end{theorem}

\paragraph{증명 전략.}
먼저 $V=\bigoplus_i V_i$와 불변성에서 ``다항식이 $T$를 소멸시키는 조건''이 각 제한연산자 소멸조건과 동치임을 보입니다. 그다음 기약인수의 지수 비교로 $m_{T_i}$의 정확한 거듭제곱 지수를 결정합니다.

\begin{proof}
정리에서 $V=\bigoplus_i V_i$를 얻었다. 각 $V_i$는 $T$-불변이므로 $T_i$가 잘 정의된다. 정의상 $V_i=\ker p_i(T)^{e_i}$이므로
\[
p_i(T_i)^{e_i}=0,
\]
따라서 $m_{T_i}\mid p_i^{e_i}$이다. $p_i$가 기약이므로
\[
m_{T_i}=p_i^{d_i}\quad(0\le d_i\le e_i)
\]
형태이다.

이제 $\operatorname{lcm}$ 식을 보이자. 다항식 $f$에 대해
\[
f(T)=0 \iff f(T_i)=0\ \text{for all }i.
\]
왼쪽에서 오른쪽은 제한으로 자명하다. 오른쪽에서 왼쪽은 임의의 $x=\sum_i x_i$ ($x_i\in V_i$)에 대해
\[
f(T)x=\sum_i f(T_i)x_i=0
\]
이므로 성립한다. 따라서 $T$를 소멸시키는 다항식의 집합은 각 $T_i$를 동시에 소멸시키는 다항식의 집합과 같고, 단위 생성원은
\[
m_T=\operatorname{lcm}(m_{T_1},\dots,m_{T_r})
\]
이다.

마지막으로 지수를 비교하자. $j\ne i$이면 $m_{T_j}\mid p_j^{e_j}$이므로 $m_{T_j}$에는 $p_i$ 인수가 없다. 따라서 위 $\operatorname{lcm}$에서 $p_i$의 지수는 오직 $m_{T_i}=p_i^{d_i}$가 제공하며 그 값은 $d_i$이다. 그런데 $m_T$에서 $p_i$의 지수는 $e_i$이므로 $d_i=e_i$이다. 결론적으로
\[
m_{T_i}=p_i^{e_i}.
\]
\end{proof}

\paragraph{의미 해석.}
프라이머리 성분은 단순히 공간을 나누는 도구가 아니라, 각 성분의 최소다항식 정보를 정확히 보존합니다. 즉 전역 최소다항식의 지수 데이터가 성분별로 분배되어 있습니다.

\begin{example}
\[
A=\operatorname{diag}\!\left(
\begin{bmatrix}1&1\\0&1\end{bmatrix},
-2,\,-2
\right)\in M_4(F)
\]
라 하자. 그러면
\[
m_A(t)=(t-1)^2(t+2)
\]
이고 프라이머리 성분은
\[
V_1=\ker(A-I)^2,\qquad V_2=\ker(A+2I)
\]
이다. Primary decomposition 정리에 의해
\[
F^4=V_1\oplus V_2,\qquad
m_{A|_{V_1}}=(t-1)^2,\quad m_{A|_{V_2}}=(t+2).
\]
\end{example}

\subsection*{주의}
\begin{warning}
프라이머리 성분은 $\ker p_i(T)$가 아니라 $\ker p_i(T)^{e_i}$로 정의됩니다. 지수 $e_i$를 생략하면 보통 분해가 완전하지 않습니다.
\end{warning}
\begin{warning}
이 직합은 일반적으로 직교직합이 아닙니다. 내적 구조를 도입하지 않은 상태에서 ``직교''라는 용어를 사용하면 안 됩니다.
\end{warning}

\subsection*{자가진단퀴즈}
\begin{enumerate}[label=\arabic*.]
\item $m_T=(t-3)^2(t^2+1)$일 때 프라이머리 성분을 정의식으로 쓰십시오.
\item $V=\bigoplus_i V_i$에서 $f(T)=0\iff f(T_i)=0\ \forall i$가 되는 이유를 한 줄로 설명하십시오.
\item Primary decomposition 증명에서 $E_i=\!b_i(T)q_i(T)$가 사영 역할을 하는 핵심 계산을 쓰십시오.
\end{enumerate}

\sectiontemplate{순환 분해 정리와 companion block}{프라이머리 분해가 ``인수별 큰 분해''라면, 순환 분해는 ``실제 계산 가능한 블록 분해''입니다. 이 절에서는 $F[t]$-모듈 관점으로 순환 블록의 존재를 보이고 companion matrix 표현까지 연결하겠습니다.}{\item $V$를 $F[t]$-모듈로 해석할 수 있습니다.\item 불변인자(invariant factor) 분해와 순환 분해 정리를 진술하고 증명할 수 있습니다.\item companion block의 블록대각 표현을 유리 표준형과 연결할 수 있습니다.}{\item PID 위 모듈\item Smith 표준형\item 순환 부분공간의 기저}

\begin{definition}
$T\in\operatorname{End}(V)$에 대해
\[
f(t)\cdot v:=f(T)v\qquad(f\in F[t],\ v\in V)
\]
로 작용을 정의하면 $V$는 $F[t]$-모듈이 된다. 이 모듈을 $V_T$로 쓴다.
\end{definition}

\begin{theorem}[유한차원 $V_T$의 불변인자 분해]
$V$가 유한차원이라 하자. 그러면 단위가 아닌 단위다항식 $a_1,\dots,a_s$가 존재하여
\[
V_T\cong F[t]/(a_1)\oplus\cdots\oplus F[t]/(a_s),\qquad
a_s\mid a_{s-1}\mid\cdots\mid a_1
\]
가 성립한다. 또한 $a_1=m_T$이다.
\end{theorem}

\paragraph{증명 전략.}
$V_T$를 유한 생성 모듈의 몫으로 제시하고, kernel을 다항식 행렬로 표현한 뒤 Smith 표준형으로 표준 분해를 얻습니다. 마지막에 모듈의 소멸 아이디얼을 비교해 가장 큰 불변인자가 최소다항식과 같음을 확인합니다.

\begin{proof}
$R=F[t]$라 두자. $V$의 $F$-기저를 $e_1,\dots,e_n$이라 하면
\[
\pi:R^n\to V_T,\qquad \pi(f_1,\dots,f_n)=\sum_{j=1}^n f_j(T)e_j
\]
는 전사 $R$-준동형이다. $K=\ker\pi$라 두면 $R$이 Noetherian PID이므로 $K$는 유한 생성이다. 생성원 $y_1,\dots,y_m\in R^n$를 잡고 열벡터 행렬
\[
A=[y_1\ \cdots\ y_m]\in M_{n,m}(R)
\]
를 만든다. 그러면 $K=\operatorname{im}(A)$이다.

PID 위 Smith 표준형 정리에 의해 가역행렬 $P\in GL_n(R)$, $Q\in GL_m(R)$와 다항식 $d_1,\dots,d_r$가 존재하여
\[
PAQ=\operatorname{diag}(d_1,\dots,d_r,0,\dots,0),\qquad d_1\mid\cdots\mid d_r.
\]
따라서
\[
R^n/K\cong R^n/\operatorname{im}(A)\cong R^n/\operatorname{im}(PAQ)
\cong \bigoplus_{i=1}^r R/(d_i)\ \oplus\ R^{\,n-r}.
\]
한편 $m_T(T)=0$이므로 임의의 $v\in V_T$는 $m_T\cdot v=0$을 만족한다. 즉 $V_T$는 torsion 모듈이며 자유부분을 가질 수 없다. 따라서 $n-r=0$이다.

단위원 $d_i$에 해당하는 항 $R/(1)$은 0모듈이므로 제거하고, 남는 다항식을 단위다항식으로 정규화하여 $a_1,\dots,a_s$라 쓰면
\[
V_T\cong \bigoplus_{j=1}^s R/(a_j),\qquad a_s\mid\cdots\mid a_1
\]
을 얻는다.

마지막으로 $a_1=m_T$를 보이자. 위 직합 모듈의 소멸 아이디얼은
\[
\bigcap_{j=1}^s (a_j)=(\operatorname{lcm}(a_1,\dots,a_s))=(a_1)
\]
이다(나눗셈 사슬 때문에 $\operatorname{lcm}=a_1$). 반면 $V_T$의 소멸 아이디얼은
\[
\{f\in R\mid f(T)=0\}=(m_T)
\]
이다. 두 아이디얼이 같으므로 $a_1=m_T$.
\end{proof}

\paragraph{의미 해석.}
유한차원 선형연산자는 $F[t]$-모듈 언어에서 ``몇 개의 순환 몫모듈의 직합''으로 완전히 분류됩니다. 이때 불변인자 사슬이 최소다항식과 블록 구조를 동시에 담습니다.

\begin{definition}
단위다항식
\[
a(t)=t^d+c_{d-1}t^{d-1}+\cdots+c_1t+c_0
\]
의 companion matrix를
\[
C(a)=
\begin{bmatrix}
0&0&\cdots&0&-c_0\\
1&0&\cdots&0&-c_1\\
0&1&\cdots&0&-c_2\\
\vdots&\vdots&\ddots&\vdots&\vdots\\
0&0&\cdots&1&-c_{d-1}
\end{bmatrix}
\]
로 정의한다.
\end{definition}

\begin{theorem}[순환 분해 정리]
위 정리의 분해
\[
V_T\cong \bigoplus_{j=1}^s F[t]/(a_j)
\]
에 대해 벡터 $v_1,\dots,v_s\in V$를 적절히 택하면
\[
V=Z(v_1;T)\oplus\cdots\oplus Z(v_s;T),\qquad
\operatorname{Ann}_T(v_j)=(a_j)
\]
가 성립한다. 따라서 적절한 기저에서 $T$의 행렬은
\[
C(a_1)\oplus\cdots\oplus C(a_s)
\]
이다.
\end{theorem}

\paragraph{증명 전략.}
모듈 동형에서 각 직합 성분의 생성원 $\overline{1}$을 끌어와 벡터 $v_j$를 정의합니다. 이어서 $f\mapsto f(T)v_j$ 사상의 kernel과 image를 계산해 각 성분이 정확히 순환 부분공간임을 보이고, 마지막에 순환 기저를 붙여 companion block 형태를 얻습니다.

\begin{proof}
$R=F[t]$라 하고, 모듈 동형
\[
\Phi:\bigoplus_{j=1}^s R/(a_j)\xrightarrow{\sim} V_T
\]
를 고정하자. $j$번째 성분의 $\overline{1}$을
\[
\mathbf{e}_j=(0,\dots,0,\overline{1},0,\dots,0)
\]
로 두고 $v_j=\Phi(\mathbf{e}_j)\in V$라 하자.

각 $j$에 대해
\[
\psi_j:R\to V,\qquad \psi_j(f)=f(T)v_j
\]
를 정의하면 $\psi_j$는 $R$-준동형이고 $\operatorname{im}\psi_j=Z(v_j;T)$이다. 또한 $\Phi^{-1}\circ\psi_j$는 $R\to R/(a_j)$의 자연사상과 같으므로
\[
\ker\psi_j=(a_j).
\]
따라서
\[
Z(v_j;T)\cong R/(a_j).
\]

한편 $\Phi$가 직합동형이므로
\[
V_T=\Phi\!\left(\bigoplus_{j=1}^s R/(a_j)\right)
=\bigoplus_{j=1}^s \Phi(R/(a_j))
=\bigoplus_{j=1}^s Z(v_j;T).
\]
즉
\[
V=Z(v_1;T)\oplus\cdots\oplus Z(v_s;T)
\]
이고 $\operatorname{Ann}_T(v_j)=(a_j)$이다.

이제 $d_j=\deg a_j$라 하자. 첫 절의 정리에 의해
\[
\mathcal{B}_j=(v_j,Tv_j,\dots,T^{d_j-1}v_j)
\]
는 $Z(v_j;T)$의 기저다. $a_j(T)v_j=0$은
\[
T^{d_j}v_j=-(c_{j,d_j-1}T^{d_j-1}+\cdots+c_{j,0}I)v_j
\]
를 주며, 따라서 $T|_{Z(v_j;T)}$의 행렬은 정확히 $C(a_j)$이다. 전체 기저 $\mathcal{B}_1\cup\cdots\cup\mathcal{B}_s$에서 $T$는 블록대각합
\[
C(a_1)\oplus\cdots\oplus C(a_s)
\]
로 표현된다.
\end{proof}

\paragraph{의미 해석.}
순환 분해는 연산자를 실제 계산 가능한 companion block들의 합으로 바꿉니다. 이 표현은 유리 표준형의 본질이며, 체가 분해체일 때 Jordan 정보로도 곧바로 연결됩니다.

\begin{example}
$\dim V=4$이고 불변인자가
\[
a_1(t)=(t-1)^2(t+1),\qquad a_2(t)=t-1
\]
이라 하자. 그러면
\[
V=Z(v_1;T)\oplus Z(v_2;T),\qquad
\operatorname{Ann}_T(v_1)=(a_1),\ \operatorname{Ann}_T(v_2)=(a_2)
\]
가 되는 벡터 $v_1,v_2$가 존재한다. 적절한 기저에서 $[T]$는
\[
C\!\big((t-1)^2(t+1)\big)\oplus C(t-1)
\]
로 주어진다. 여기서 최소다항식은 $a_1$, 특성다항식은 $a_1a_2$이다.
\end{example}

\subsection*{주의}
\begin{warning}
순환 분해에서 생성벡터 $v_j$는 유일하지 않습니다. 유일한 것은 다항식 사슬 $a_s\mid\cdots\mid a_1$입니다.
\end{warning}
\begin{warning}
프라이머리 분해와 순환 분해는 서로 다른 질문에 답합니다. 전자는 ``기약인수별 분리'', 후자는 ``계산 블록의 표준화''를 담당합니다.
\end{warning}

\subsection*{자가진단퀴즈}
\begin{enumerate}[label=\arabic*.]
\item $V_T\cong F[t]/(a_1)\oplus F[t]/(a_2)$이고 $a_2\mid a_1$일 때 왜 $m_T=a_1$인지 설명하십시오.
\item $a(t)=t^3-2t+1$의 companion matrix를 직접 쓰십시오.
\item 순환 분해에서 특성다항식이 $\prod_j a_j$가 되는 이유를 차수 관점으로 설명하십시오.
\end{enumerate}

\chapterapplicationtemplate{상태전이행렬 $A\in M_n(F)$가 주어졌을 때, 먼저 $m_A$를 인수분해하여 프라이머리 성분을 분리한 뒤 각 성분에서 순환 분해를 수행해 block companion 형태를 구성해 보십시오. 이 표현을 사용하면 $A^k\mathbf{x}_0$ 계산, 선형 점화식의 최소 차수 판정, 그리고 Jordan 형식 계산 전 단계의 데이터 정리가 체계적으로 이루어집니다.}

\chaptersummarytemplate

\section*{종합 연습문제}

\subsection*{연습문제 A (기초 확인)}
\begin{enumerate}[label=A\arabic*.]
\item $N=\begin{bmatrix}0&1&0\\0&0&1\\0&0&0\end{bmatrix}$, $v=e_3$에 대해 $\operatorname{Ann}_N(v)$와 $Z(v;N)$의 기저를 구하십시오.
\item $m_T(t)=(t-2)^3(t+1)^2$일 때 프라이머리 성분을 정의식으로 쓰십시오.
\item 서로소인 $f,g\in F[t]$에 대해 $\ker((fg)(T))=\ker f(T)\oplus\ker g(T)$ 공식을 $T=\operatorname{diag}(1,1,-2)$, $f=t-1$, $g=t+2$에서 확인하십시오.
\item
\[
A=\operatorname{diag}\!\left(\begin{bmatrix}1&1\\0&1\end{bmatrix},-2,-2\right)
\]
에 대해 $V_1=\ker(A-I)^2$, $V_2=\ker(A+2I)$의 차원을 구하십시오.
\item $a(t)=t^4+2t^2+1$의 companion matrix $C(a)$를 쓰십시오.
\item $\deg m_{T,v}=5$이면 $\dim Z(v;T)=5$임을 설명하십시오.
\item $V=Z(v;T)$이고 $m_{T,v}=t^3-t$일 때 $T^3v$를 $\{v,Tv,T^2v\}$의 선형결합으로 표현하십시오.
\item $V=V_1\oplus V_2$에서 $f(T)|_{V_1}=0$, $f(T)|_{V_2}=0$이면 $f(T)=0$임을 보이십시오.
\end{enumerate}

\subsection*{연습문제 B (개념 연결)}
\begin{enumerate}[label=B\arabic*.]
\item 부분집합 $S_1,S_2\subseteq V$에 대해
\[
\operatorname{Ann}_T(S_1\cup S_2)=\operatorname{Ann}_T(S_1)\cap\operatorname{Ann}_T(S_2)
\]
를 증명하십시오.
\item Primary decomposition 증명에서 정의한 $E_i=b_i(T)q_i(T)$가
\[
E_i^2=E_i,\qquad E_iE_j=0\ (i\ne j),\qquad \sum_i E_i=I
\]
를 만족함을 보이십시오.
\item $V=\bigoplus_{j=1}^s Z(v_j;T)$, $\operatorname{Ann}_T(v_j)=(a_j)$, $a_s\mid\cdots\mid a_1$일 때 $m_T=a_1$임을 증명하십시오.
\item $T$의 불변인자가 $a_1,a_2,\dots,a_s$일 때 $\chi_T(t)=\prod_{j=1}^s a_j(t)$임을 보이십시오.
\item $m_T$가 중근이 없는 다항식이면(분해체에서) $T$가 대각화 가능함을 설명하십시오.
\end{enumerate}

\subsection*{연습문제 C (증명/종합)}
\begin{enumerate}[label=C\arabic*.]
\item $m_T=\prod_{i=1}^r p_i^{e_i}$일 때
\[
V_i=\ker p_i(T)^{e_i}
\]
가 $T$에 의해 유일하게 결정되는 성분임을 증명하십시오.
\item 불변인자 분해
\[
V_T\cong\bigoplus_{j=1}^s F[t]/(a_j),\qquad a_s\mid\cdots\mid a_1
\]
에서 다항식 사슬 $\{a_j\}$가 순서를 제외하면 유일함을 증명하십시오.
\item $F=\C$에서 순환 분해와 Jordan 분해가 어떻게 호환되는지 증명하십시오. 즉 각 $a_j$의 선형인수 거듭제곱 분해가 Jordan 블록 크기 정보로 바뀌는 과정을 엄밀히 서술하십시오.
\end{enumerate}